As described in Chapter~\ref{chapter:Gravitational_wave_interferometers}, gravitational wave detections are made with extremely sensitive interferometers. As the LIGO interferometers cannot meet their design sensitivity with their current hardware, upgrades to the detectors will be made. The two upgrades explored in this thesis are: the high-powered laser for LIGO's fourth observation run and the balanced homodyne detector for LIGO A+.
\par
A high laser power is needed to reduce the shot noise limit of an interferometer, so the power that the LIGO detectors are designed to use is 125\,W. The original concept for LIGO's laser centred around the HPO; this was a ring cavity containing laser crystals which was seeded with a low noise laser, and it could produce $\sim150\,\si{\watt}$ of power in the HG00 mode. Unfortunately, it was discovered that this design of laser coupled pointing noise to the differential arm signal. This noise originated from the HPO's water-cooling system, so the HPO was replaced with a single-pass amplifier which exhibited less pointing noise. This amplifier was used throughout LIGO's third observation run; however, the maximum power of this laser was limited to $\sim70\si{\watt}$.
\par
To increase the maximum power that can be delivered to the interferometers for LIGO's fourth observation run, two single-pass amplifiers will be used in series. In Chapter~\ref{chapter:TnT_lab_work}, the prototype for this laser was characterised. With this laser, $\sim100\,\si{\watt}$ of amplitude stabilised light in the HG00 mode was generated. As the pointing noise of this prototype was shown to exceed LIGO's requirement, sources of pointing noise, such as turbulence in the laser's water-cooling system, must be kept minimal when it is commissioned to be part of the interferometer.
\par
Two actuators for stabilising the amplitude of the prototype laser were investigated. These were an AOM and a current shunt on one of the pump diodes in one of the amplifiers. Because the operators of the interferometers will need to have control over the input laser power, the AOM was chosen to be the actuator used to stabilise the amplitude of the laser as it was found to behave consistently over a wider range of power. 
\par
Balanced homodyne detection is a crucial part of the upgrade from advanced LIGO to LIGO A+ since the homodyne angle must be a free parameter if the detectors are to reach the planned quantum noise by using squeezed light. The fundamentals of balanced homodyne detection and the phase noise requirement for the local oscillator were discussed in Chapter~\ref{chapter:BHD_for_ligo_1}. 
\par
To preserve a squeezed state of light, it is essential that there is minimal loss within the detector. As part of the balanced homodyne detector upgrade, active wavefront control will be used to minimise the loss due to mode mismatch between the interferometer's arm mode and the output mode cleaners. To tackle this problem, the range of interferometer arm modes that the active wavefront control will have to mode match was determined by simulation. This is described in Chapter~\ref{chapter:simulation_finesse}. The dominant source of uncertainty in the mode emerging from the SRM comes from the uncertainty in the radius of curvatures for the optics within the SRC. The probability of a mode exiting the SRM was determined by the probability of the optics within the SRC having the radii of curvature required to produce the beam. Combining this with measurements of the Gouy phase of the LLO SRC, the beam exiting the SRM was inferred to be 1.84\,mm in width and to have a defocus of $-2.80$ Dioptre.
\par
With a set of likely modes that the adaptive optics for mode matching would need to correct for, the robustness of the active wavefront control for the balanced homodyne detector to a mode mismatch was determined. This is explored in Chapter~\ref{chapter:BHD_for_ligo_2}. With the pessimistic assumption that the range of the active optics will be $\pm50\,\rm{mD}$, it was found that most of the likely beams that could emerge from the SRM are covered by the proposed design for the active wavefront control. 
\par
Finally, we looked to the future of ground-based gravitational wave detection. As the thermal noise of the current test masses' coatings and substrates will limit the sensitivity of the A+ detectors, and because their noise properties do not improve when they are cryogenically cooled, research into new test mass technology is underway. There have been promising results in this area of research; however, the new test masses would require the interferometers to be operated with a laser of either 1.5\,\si{\micro\metre} or 2\,\si{\micro\metre} in wavelength. Unlike for 1.5\,\si{\micro\metre} wavelength light, many technologies for 2\,\si{\micro\metre} light that the interferometers rely on are not mature enough for the interferometers to reach their design sensitivities.
\par
Photodiodes are a key element of an interferometer, so off-the-shelf extended InGaAs photodiodes for detecting 2\,\si{\micro\metre} light were characterised in Chapter~\ref{chapter:2um_photoiode_chapter}. It is crucial that the quantum efficiency of a photodiode is at least 99\% so that future gravitational wave detectors can reach the degree of squeezing required to meet their design sensitivity. However, the strained nature of extended InGaAs means that making high quantum efficiency photodiodes is challenging. 
\par
It was found that off-the-shelf extended InGaAs photodiodes do not meet the requirements of a third-generation ground-based interferometer under any operating conditions because their quantum efficiencies are too low. The quantum efficiency of extended InGaAs photodiodes can be improved by reverse biasing them; however, when reverse biased, extended InGaAs photodiodes will exhibit $1/f$ dark noise. 
\par
The $1/f$ noise arises due to the generation-recombination current associated with the strain-induced by defects in the extended InGaAs photodiode. The dark noise for a selection of extended InGaAs photodiodes as a function of bias and temperature was shown in Chapter~\ref{chapter:2um_photoiode_chapter}. If biased with the maximum stated value in their datasheets, the $1/f$ produced by the photodiodes would exceed the typical amount of shot noise of the light detected by the photodiode. As expected, cooling the photodiodes caused the dark noise to be exponentially reduced.
\par
As the quantum efficiency and dark noise problems surrounding the use of extended InGaAs photodiodes are due to the strain induced by the lattice mismatch and the defects caused by this, research into the design and the optimal growth conditions of extended InGaAs is needed. 